%%%%%%%%%%  scan idx 18  |  folio 5  %%%%%%%%%%
est un
isomorphisme modulo torsion, ce qui justifie le point de vue suivant lequel cet anneau tensorisé par $\mathbb Q$ peut jouer en tous points le même rôle que l'anneau de Chow tensorisé par $\mathbb Q$. D'autre part, on a fait exactement ce qu'il fallait pour que la formule (2.3) (où on remplace $\operatorname{Filt}^i_{\mathrm{top}}$ par $\operatorname{Filt}^i$) définisse une théorie des classes de Chern à valeurs dans $\Gr^\bullet(X)$, dont les classes de Chern dans $\Gr^{\bullet}_{\mathrm{top}}(X)$ se déduisent à l'aide de l'homomorphisme (2.4). On montre (Exp.\ V), grâce au fait (Exp.\ VI) que $K(X)$ est un $\lambda$-anneau \underline{spécial} au sens de [RRR] (i.e.\ un $\lambda$-anneau au sens de Exp.\ V) que cette notion de classe de Chern satisfait à toutes les propriétés formelles habituelles, passées en revue dans [RRR].

\subsubsection*{2.4.}
On notera un avantage très appréciable de la définition de $\Gr^\bullet(X)$ sur celle de $A(X)$ et $\Gr^{\bullet}_{\mathrm{top}}(X)$, c'est qu'elle garde un sens pour un schéma $X$ absolument quelconque, et plus généralement pour un topos localement annelé quelconque $X$, en prenant alors pour $K(X)$ le groupe de classes construit avec les Modules localement libres sur $X$, qui est encore un $\lambda$-anneau augmenté, donc muni d'une filtration correspondante, permettant de définir un gradué associé $\Gr^\bullet(X)$. Comme prix de cette généralité, il faut cependant remarquer que lorsqu'on n'est pas disposé à négliger les phénomènes de torsion i.e.\ à tensoriser par $\mathbb Q$, les propriétés de $\Gr^\bullet(X)$ sont moins satisfaisantes à divers égards que celles de ses deux concurents (Exp.\ XIV n\textsuperscript{o}~4). C'est ainsi que, supposant à nouveau que $f:X\to Y$ est un morphisme propre de schémas noethériens réguliers munis de Modules inversibles amples, il n'est pas vrai en général que l'homomorphisme $f_*:K(X)\longrightarrow K(Y)$ applique $\operatorname{Filt}^i K(X)$ dans $\operatorname{Filt}^{i-d} K(Y)$, où $d$ est la dimension relative de $X$ sur $Y$, et définisse par suite par passage aux gradués associés un homomorphisme de degré $-d$ de $\Gr^\bullet(X)$ dans $\Gr^\bullet(Y)$. Cependant, on peut prouver (Exp.\ VII) qu'il en est ainsi à condition de tensoriser par $\mathbb Q$, de sorte qu'on obtient un homomorphisme de degré $-d$
\begin{equation}\tag{2.5}
f_*:\Gr^\bullet(X)_{\mathbb Q}\longrightarrow \Gr^\bullet(Y)_{\mathbb Q},
\end{equation}
jouant le rôle de l'homomorphisme «~image directe de cycles~» $A(X)\to A(Y)$.

Ceci posé, nous voyons donc que nous sommes en possession de tous les éléments de structure pour donner un sens à la formule (1.2) lorsqu'on ne postule pas l'existence d'un corps de base, $f:X\to Y$ étant morphisme propre de schémas noethériens réguliers admettant des Modules inversibles amples.

